% common.tex -- everything shared by a4.tex and letter.tex.
% Defines the packages, fonts, suit symbols, the two card faces (\sideA = Rules,
% \sideB = Scoring), the 4-up and 1-up layouts, and \euchrerules which typesets
% both pages. Define \oneupmode before \input-ing the paper file to get the 1-up
% layout (one card per page, the 4-up cell scaled 2x), e.g.
%   pdflatex -jobname=euchre-rules-1up '\def\oneupmode{}% Bid Euchre reference cards, 4-up double-sided, US Letter.
% Build: pdflatex letter.tex   (run twice so the cut lines settle)
\documentclass{article}
\usepackage[letterpaper,margin=0pt]{geometry}
% common.tex -- everything shared by a4.tex and letter.tex.
% Defines the packages, fonts, suit symbols, the two card faces (\sideA = Rules,
% \sideB = Scoring), the 4-up and 1-up layouts, and \euchrerules which typesets
% both pages. Define \oneupmode before \input-ing the paper file to get the 1-up
% layout (one card per page, the 4-up cell scaled 2x), e.g.
%   pdflatex -jobname=euchre-rules-1up '\def\oneupmode{}% Bid Euchre reference cards, 4-up double-sided, US Letter.
% Build: pdflatex letter.tex   (run twice so the cut lines settle)
\documentclass{article}
\usepackage[letterpaper,margin=0pt]{geometry}
\input{common}

\begin{document}
\euchrerules
\end{document}
'
%
% It expects the including file to have already loaded \documentclass and geometry
% (zero margins), because the paper size is the only thing that differs.
%
% Print double-sided (flip on LONG edge) at 100% / "Actual size", then cut along the
% dashed centre lines (printed on the front page only). All four cells on a page are
% identical, so left/right mirroring is not an issue.
%
% Bid Euchre with a 28-card deck (A to 8), 7 cards each. One round of bidding: 1 to 7
% tricks plus a trump suit, high (no trump, ace high) or low (no trump, 8 high); the
% dealer must bid at least 4 if everyone passes. On a contract of 5, 6 or 7 the maker
% may call salt (both partnerships swap one card between partners), pepper (maker plays
% alone) or salt & pepper (swap first, then play alone). See \sideB for the scoring.

\usepackage[scaled]{helvet}
\usepackage[T1]{fontenc}
\usepackage{array}
\usepackage{tabularx}
\usepackage{xcolor}
\usepackage{pifont}
\usepackage{enumitem}
\usepackage{url}\urlstyle{sf}
\usepackage{tikz}
\renewcommand{\familydefault}{\sfdefault}
\pagestyle{empty}
\setlength{\parindent}{0pt}
\setlength{\parskip}{0pt}
\hyphenpenalty=10000 \exhyphenpenalty=10000 \tolerance=2000 \emergencystretch=2em

% ---- CONFIGURE ----
\newcommand{\footurl}{https://github.com/michal-dobrogost/minibridge/tree/main/euchre}  % footnote on the scoring card
\newlength{\inset}\setlength{\inset}{0.25in}   % white margin inside each cell
\newcommand{\bodysize}{\fontsize{9}{10.8}\selectfont}
\newcommand{\titlesize}{\fontsize{19}{21}\selectfont}
\newcommand{\gap}{\vspace{0.7\baselineskip}}
\newcommand{\winscore}{33}      % points needed to win the game
% -------------------

% table header cell split over two lines
\newcommand{\stackhead}[1]{{\renewcommand{\arraystretch}{1.0}\begin{tabular}[t]{@{}l@{}}#1\end{tabular}}}

% suit symbols
\newcommand{\Club}{\ding{168}}
\newcommand{\Diamondsuit}{{\color{red}\ding{169}}}
\newcommand{\Heart}{{\color{red}\ding{170}}}
\newcommand{\Spade}{\ding{171}}

% 2x2 block of suits, then the title to its right
\newcommand{\cardtitle}[1]{%
  \noindent\hbox to \linewidth{%
    \vbox{\fontsize{15}{16}\selectfont\offinterlineskip
      \hbox to 2.1em{\hss\Spade\hss\Heart\hss}\vskip 0.15em
      \hbox to 2.1em{\hss\Diamondsuit\hss\Club\hss}}%
    \hspace{0.5em}%
    \vbox{\hsize=\dimexpr\linewidth-3em\relax\titlesize\bfseries #1}\hss}%
  \par\vspace{1.1\baselineskip}}

\newlength{\cellw}\newlength{\cellh}
\setlength{\cellw}{0.5\paperwidth}
\setlength{\cellh}{0.5\paperheight}

% ---------- SIDE A: RULES ----------
\newcommand{\sideA}{%
  \cardtitle{Toronto Bid Euchre: Rules}
  {\raggedright\begin{enumerate}[leftmargin=1.5em,itemsep=0.5\baselineskip,topsep=0pt,parsep=0pt]
    \item Four players form two teams with partners sitting across from each other.
      Use a 28-card deck: A, K, Q, J, 10, 9, 8 in each suit.
    \item Deal the cards, clockwise, starting left of dealer. Seven cards each.
    \item Each player bids once, starting left of dealer.
      Each player passes or bids a number of tricks their side is to take (1--7) plus a trump suit, high or low
      (e.g.\ 4\Heart, 3~high, 5~low).
      Each bid must be for more tricks than the previous bid.
      If everyone else passes, the dealer must bid at least 4.
      The highest bidder is known as the maker.
    \item
      \emph{(Beginners skip)} On a contract of 5, 6 or 7 tricks, the maker may call:\\
      Salt --- all players swap one card face down with their partner.\\
      Pepper --- maker plays alone (partner sits out with hand face down).\\
      Salt \& pepper --- salt first, then pepper.
    \item
      Maker leads one card face up for the first trick; others follow clockwise.
      \emph{You must follow the suit led.} If you can't, play trump or any other card.
      The highest trump wins the trick; otherwise the highest card of the suit led wins.
      The winner leads to the next trick.
    \item
      Card ranking changes based on the winning bid (the contract).
      In a trump contract the jack of trump (right bower) is highest, the other jack of the same colour (left bower) is second highest.
      \\[0.3\baselineskip]
      {\renewcommand{\arraystretch}{1.2}%
      \begin{tabularx}{\linewidth}{|l|>{\raggedright\arraybackslash}X|}
        \hline
        Contract & Ranking, highest first \\
        \hline
        \stackhead{Trump suit (e.g.\ \Heart)} &
          (E.g. J\Heart, J\Diamondsuit, A\Heart, K\Heart, Q\Heart, 10\Heart, 9\Heart, 8\Heart) \\
        \hline
        High (no trump) & A, K, Q, J, 10, 9, 8 \\
        \hline
        Low (no trump) & 8, 9, 10, J, Q, K, A \\
        \hline
      \end{tabularx}}\\[0.3\baselineskip]
      \emph{The left bower (e.g. J\Diamondsuit) is a trump, not a card of its own suit.}
    \item Tally score. For the next hand, dealer is the next player clockwise.
  \end{enumerate}}
}

% ---------- SIDE B: SCORING ----------
\newcommand{\sideB}{%
  \raggedright
  \renewcommand{\gap}{\vspace{0.5\baselineskip}}% tighter than the rules side, so it all fits
  \cardtitle{Toronto Bid Euchre: Score}

  \textbf{Makers}\gap

  The maker and partner are the makers. The other side are the defenders.
  The contract is made if the makers win at least the number of tricks bid, otherwise it is set.
  Makers score points as follows:
  \gap

  \renewcommand{\arraystretch}{1.3}
  \begin{tabular}{|l|l|l|}
    \hline
    Contract & Made & Set (failed) \\
    \hline
    Plain & tricks won & \textminus{}tricks bid \\
    Salted & tricks won \textminus\ 1 & \textminus{}tricks bid \textminus\ 1 \\
    Peppered & tricks won + 4 & \textminus{}tricks bid \textminus\ 4 \\
    Salted \& peppered & tricks won + 2 & \textminus{}tricks bid \textminus\ 8 \\
    \hline
  \end{tabular}\gap

  \textbf{Defenders}\gap

  The defenders score 1 point for each trick they take, whether or not the contract is made.\gap

  \textbf{Examples}\gap

  \begin{tabular}{|l|l|r@{}l|l|}
    \hline
    Contract & \stackhead{Tricks won\\by makers} & \multicolumn{2}{l|}{\stackhead{Makers\\score}} & \stackhead{Defenders\\score} \\
    \hline
    Plain 6 low                & 6 & 6            &                                          & 1 \\
    Salted 5\Heart             & 6 & 5            & \ =\ \phantom{\textminus}6 \textminus\ 1 & 1 \\
    Peppered 5\Heart           & 5 & 9            & \ =\ \phantom{\textminus}5 + 4           & 2 \\
    Salted \& peppered 5\Heart & 4 & \textminus13 & \ =\ \textminus5 \textminus\ 8           & 3 \\
    \hline
  \end{tabular}\gap

  \textbf{Winning}\gap

  %The first side to reach \winscore{} points wins (about 10 hands). If both teams achieve 33 points on the same hand, the higher score count wins (if still tied then makers on last deal win). Scores can go below zero.
  First side to \winscore{} wins (about 10 hands). If both get there on the same hand, the higher score wins and makers take ties. Scores can go below zero.

  \gap
  \vfill{\fontsize{5.5}{6.5}\selectfont\noindent\expandafter\url\expandafter{\footurl}}}

% ---------- layout ----------
\newcommand{\cell}[1]{%
  \vbox to \cellh{\vss\hbox to \cellw{\hss
    \parbox[c][\dimexpr\cellh-2\inset\relax][t]{\dimexpr\cellw-2\inset\relax}{\bodysize #1}%
  \hss}\vss}}
\newcommand{\cellrow}[1]{\hbox to \paperwidth{\cell{#1}\cell{#1}}}

\newcommand{\fourup}[1]{%
  \noindent\vbox to \paperheight{\offinterlineskip\cellrow{#1}\cellrow{#1}}}

% Dashed cut lines through the page centre; drawn on the front page only, so the
% back side prints clean. Must follow \fourup in the same paragraph.
\newcommand{\cutlines}{%
  \begin{tikzpicture}[remember picture,overlay]
    \draw[gray!60,dash pattern=on 3pt off 3pt,thin]
      (current page.north) -- (current page.south);
    \draw[gray!60,dash pattern=on 3pt off 3pt,thin]
      (current page.west) -- (current page.east);
  \end{tikzpicture}}

% One cell (laid out exactly as in the 4-up version) scaled 2x to fill the page.
\newcommand{\oneup}[1]{%
  \noindent\vbox to \paperheight{\offinterlineskip\scalebox{2}{\cell{#1}}\vss}}

% Page 1 = side A (Rules), page 2 = side B (Scoring); four copies each in 4-up mode.
\ifdefined\oneupmode
  \newcommand{\euchrerules}{%
    \oneup{\sideA}
    \newpage
    \oneup{\sideB}
  }
\else
  \newcommand{\euchrerules}{%
    \fourup{\sideA}\cutlines
    \newpage
    \fourup{\sideB}
  }
\fi


\begin{document}
\euchrerules
\end{document}
'
%
% It expects the including file to have already loaded \documentclass and geometry
% (zero margins), because the paper size is the only thing that differs.
%
% Print double-sided (flip on LONG edge) at 100% / "Actual size", then cut along the
% dashed centre lines (printed on the front page only). All four cells on a page are
% identical, so left/right mirroring is not an issue.
%
% Bid Euchre with a 28-card deck (A to 8), 7 cards each. One round of bidding: 1 to 7
% tricks plus a trump suit, high (no trump, ace high) or low (no trump, 8 high); the
% dealer must bid at least 4 if everyone passes. On a contract of 5, 6 or 7 the maker
% may call salt (both partnerships swap one card between partners), pepper (maker plays
% alone) or salt & pepper (swap first, then play alone). See \sideB for the scoring.

\usepackage[scaled]{helvet}
\usepackage[T1]{fontenc}
\usepackage{array}
\usepackage{tabularx}
\usepackage{xcolor}
\usepackage{pifont}
\usepackage{enumitem}
\usepackage{url}\urlstyle{sf}
\usepackage{tikz}
\renewcommand{\familydefault}{\sfdefault}
\pagestyle{empty}
\setlength{\parindent}{0pt}
\setlength{\parskip}{0pt}
\hyphenpenalty=10000 \exhyphenpenalty=10000 \tolerance=2000 \emergencystretch=2em

% ---- CONFIGURE ----
\newcommand{\footurl}{https://github.com/michal-dobrogost/minibridge/tree/main/euchre}  % footnote on the scoring card
\newlength{\inset}\setlength{\inset}{0.25in}   % white margin inside each cell
\newcommand{\bodysize}{\fontsize{9}{10.8}\selectfont}
\newcommand{\titlesize}{\fontsize{19}{21}\selectfont}
\newcommand{\gap}{\vspace{0.7\baselineskip}}
\newcommand{\winscore}{33}      % points needed to win the game
% -------------------

% table header cell split over two lines
\newcommand{\stackhead}[1]{{\renewcommand{\arraystretch}{1.0}\begin{tabular}[t]{@{}l@{}}#1\end{tabular}}}

% suit symbols
\newcommand{\Club}{\ding{168}}
\newcommand{\Diamondsuit}{{\color{red}\ding{169}}}
\newcommand{\Heart}{{\color{red}\ding{170}}}
\newcommand{\Spade}{\ding{171}}

% 2x2 block of suits, then the title to its right
\newcommand{\cardtitle}[1]{%
  \noindent\hbox to \linewidth{%
    \vbox{\fontsize{15}{16}\selectfont\offinterlineskip
      \hbox to 2.1em{\hss\Spade\hss\Heart\hss}\vskip 0.15em
      \hbox to 2.1em{\hss\Diamondsuit\hss\Club\hss}}%
    \hspace{0.5em}%
    \vbox{\hsize=\dimexpr\linewidth-3em\relax\titlesize\bfseries #1}\hss}%
  \par\vspace{1.1\baselineskip}}

\newlength{\cellw}\newlength{\cellh}
\setlength{\cellw}{0.5\paperwidth}
\setlength{\cellh}{0.5\paperheight}

% ---------- SIDE A: RULES ----------
\newcommand{\sideA}{%
  \cardtitle{Toronto Bid Euchre: Rules}
  {\raggedright\begin{enumerate}[leftmargin=1.5em,itemsep=0.5\baselineskip,topsep=0pt,parsep=0pt]
    \item Four players form two teams with partners sitting across from each other.
      Use a 28-card deck: A, K, Q, J, 10, 9, 8 in each suit.
    \item Deal the cards, clockwise, starting left of dealer. Seven cards each.
    \item Each player bids once, starting left of dealer.
      Each player passes or bids a number of tricks their side is to take (1--7) plus a trump suit, high or low
      (e.g.\ 4\Heart, 3~high, 5~low).
      Each bid must be for more tricks than the previous bid.
      If everyone else passes, the dealer must bid at least 4.
      The highest bidder is known as the maker.
    \item
      \emph{(Beginners skip)} On a contract of 5, 6 or 7 tricks, the maker may call:\\
      Salt --- all players swap one card face down with their partner.\\
      Pepper --- maker plays alone (partner sits out with hand face down).\\
      Salt \& pepper --- salt first, then pepper.
    \item
      Maker leads one card face up for the first trick; others follow clockwise.
      \emph{You must follow the suit led.} If you can't, play trump or any other card.
      The highest trump wins the trick; otherwise the highest card of the suit led wins.
      The winner leads to the next trick.
    \item
      Card ranking changes based on the winning bid (the contract).
      In a trump contract the jack of trump (right bower) is highest, the other jack of the same colour (left bower) is second highest.
      \\[0.3\baselineskip]
      {\renewcommand{\arraystretch}{1.2}%
      \begin{tabularx}{\linewidth}{|l|>{\raggedright\arraybackslash}X|}
        \hline
        Contract & Ranking, highest first \\
        \hline
        \stackhead{Trump suit (e.g.\ \Heart)} &
          (E.g. J\Heart, J\Diamondsuit, A\Heart, K\Heart, Q\Heart, 10\Heart, 9\Heart, 8\Heart) \\
        \hline
        High (no trump) & A, K, Q, J, 10, 9, 8 \\
        \hline
        Low (no trump) & 8, 9, 10, J, Q, K, A \\
        \hline
      \end{tabularx}}\\[0.3\baselineskip]
      \emph{The left bower (e.g. J\Diamondsuit) is a trump, not a card of its own suit.}
    \item Tally score. For the next hand, dealer is the next player clockwise.
  \end{enumerate}}
}

% ---------- SIDE B: SCORING ----------
\newcommand{\sideB}{%
  \raggedright
  \renewcommand{\gap}{\vspace{0.5\baselineskip}}% tighter than the rules side, so it all fits
  \cardtitle{Toronto Bid Euchre: Score}

  \textbf{Makers}\gap

  The maker and partner are the makers. The other side are the defenders.
  The contract is made if the makers win at least the number of tricks bid, otherwise it is set.
  Makers score points as follows:
  \gap

  \renewcommand{\arraystretch}{1.3}
  \begin{tabular}{|l|l|l|}
    \hline
    Contract & Made & Set (failed) \\
    \hline
    Plain & tricks won & \textminus{}tricks bid \\
    Salted & tricks won \textminus\ 1 & \textminus{}tricks bid \textminus\ 1 \\
    Peppered & tricks won + 4 & \textminus{}tricks bid \textminus\ 4 \\
    Salted \& peppered & tricks won + 2 & \textminus{}tricks bid \textminus\ 8 \\
    \hline
  \end{tabular}\gap

  \textbf{Defenders}\gap

  The defenders score 1 point for each trick they take, whether or not the contract is made.\gap

  \textbf{Examples}\gap

  \begin{tabular}{|l|l|r@{}l|l|}
    \hline
    Contract & \stackhead{Tricks won\\by makers} & \multicolumn{2}{l|}{\stackhead{Makers\\score}} & \stackhead{Defenders\\score} \\
    \hline
    Plain 6 low                & 6 & 6            &                                          & 1 \\
    Salted 5\Heart             & 6 & 5            & \ =\ \phantom{\textminus}6 \textminus\ 1 & 1 \\
    Peppered 5\Heart           & 5 & 9            & \ =\ \phantom{\textminus}5 + 4           & 2 \\
    Salted \& peppered 5\Heart & 4 & \textminus13 & \ =\ \textminus5 \textminus\ 8           & 3 \\
    \hline
  \end{tabular}\gap

  \textbf{Winning}\gap

  %The first side to reach \winscore{} points wins (about 10 hands). If both teams achieve 33 points on the same hand, the higher score count wins (if still tied then makers on last deal win). Scores can go below zero.
  First side to \winscore{} wins (about 10 hands). If both get there on the same hand, the higher score wins and makers take ties. Scores can go below zero.

  \gap
  \vfill{\fontsize{5.5}{6.5}\selectfont\noindent\expandafter\url\expandafter{\footurl}}}

% ---------- layout ----------
\newcommand{\cell}[1]{%
  \vbox to \cellh{\vss\hbox to \cellw{\hss
    \parbox[c][\dimexpr\cellh-2\inset\relax][t]{\dimexpr\cellw-2\inset\relax}{\bodysize #1}%
  \hss}\vss}}
\newcommand{\cellrow}[1]{\hbox to \paperwidth{\cell{#1}\cell{#1}}}

\newcommand{\fourup}[1]{%
  \noindent\vbox to \paperheight{\offinterlineskip\cellrow{#1}\cellrow{#1}}}

% Dashed cut lines through the page centre; drawn on the front page only, so the
% back side prints clean. Must follow \fourup in the same paragraph.
\newcommand{\cutlines}{%
  \begin{tikzpicture}[remember picture,overlay]
    \draw[gray!60,dash pattern=on 3pt off 3pt,thin]
      (current page.north) -- (current page.south);
    \draw[gray!60,dash pattern=on 3pt off 3pt,thin]
      (current page.west) -- (current page.east);
  \end{tikzpicture}}

% One cell (laid out exactly as in the 4-up version) scaled 2x to fill the page.
\newcommand{\oneup}[1]{%
  \noindent\vbox to \paperheight{\offinterlineskip\scalebox{2}{\cell{#1}}\vss}}

% Page 1 = side A (Rules), page 2 = side B (Scoring); four copies each in 4-up mode.
\ifdefined\oneupmode
  \newcommand{\euchrerules}{%
    \oneup{\sideA}
    \newpage
    \oneup{\sideB}
  }
\else
  \newcommand{\euchrerules}{%
    \fourup{\sideA}\cutlines
    \newpage
    \fourup{\sideB}
  }
\fi
